\documentclass[11pt,a4paper]{article}

\usepackage[T1]{fontenc}
\usepackage[utf8]{inputenc}
\usepackage{geometry}
\usepackage{hyperref}
\usepackage{longtable}
\usepackage{array}
\usepackage{enumitem}
\usepackage{xcolor}
\usepackage{listings}
\usepackage{titlesec}
\usepackage{url}
\usepackage{xurl}
\usepackage{microtype}
\usepackage{ragged2e}

\geometry{margin=2.1cm}
\hypersetup{
  colorlinks=true,
  linkcolor=blue!55!black,
  urlcolor=blue!55!black
}
\urlstyle{same}
\setlist{nosep}
\setlength{\parskip}{0.45em}
\setlength{\parindent}{0pt}
\setlength{\tabcolsep}{4pt}
\renewcommand{\arraystretch}{1.15}
\emergencystretch=3em
\sloppy
\titleformat{\section}{\Large\bfseries}{\thesection}{0.6em}{}
\titleformat{\subsection}{\large\bfseries}{\thesubsection}{0.6em}{}
\titleformat{\subsubsection}{\normalsize\bfseries}{\thesubsubsection}{0.6em}{}

\lstdefinestyle{pythonapi}{
  basicstyle=\ttfamily\small,
  breaklines=true,
  columns=fullflexible,
  frame=single,
  backgroundcolor=\color{black!3}
}

\newcommand{\codepath}[1]{\nolinkurl{#1}}
\newcommand{\codeheading}[1]{\texorpdfstring{\codepath{#1}}{#1}}
\newcommand{\codesym}[1]{\nolinkurl{#1}}
\newcommand{\apiitem}[2]{\item \codesym{#1}: #2}
\newcommand{\apisym}[1]{\item \codesym{#1}}
\newcolumntype{L}[1]{>{\RaggedRight\arraybackslash}p{#1}}
\newcolumntype{P}[1]{>{\RaggedRight\ttfamily\arraybackslash}p{#1}}

\title{GPS L1 C/A Offline Decoder and Visualizer\\Code Documentation}
\author{Fraunhofer FHR Educational GPS Decoder Repository}
\date{Generated for repository state on 2026-05-06}

\begin{document}
\maketitle
\tableofcontents
\newpage

\section{Scope}

This document describes all source code in the repository rooted at \codepath{Fraunhofer_FHR}. It covers the Python application under \codepath{app/}, the pytest suite under \codepath{app/tests/}, and the PowerShell repository tools under \codepath{tools/}. Large local IQ captures such as \codepath{*.bin} files are intentionally excluded because they are analysis assets, not source code.

The application is an offline, didactic GPS L1 C/A decoder and visualizer. Its code is organized around a deliberately transparent signal-processing pipeline:

\begin{enumerate}
  \item inspect and load little-endian \codepath{complex64} IQ recordings,
  \item acquire PRN-specific C/A signals over Doppler and code phase,
  \item track each acquired satellite with simple carrier and code loops,
  \item form 20 ms navigation-bit decisions from prompt integrations,
  \item frame and verify GPS LNAV words and subframes,
  \item decode ephemeris evidence where available,
  \item solve an offline PVT/time result from decoded navigation and pseudorange evidence,
  \item present each stage in a Qt GUI with PRN-specific diagnostics.
\end{enumerate}

\section{Repository Layout}

\begin{longtable}{P{0.30\linewidth}L{0.64\linewidth}}
\textbf{Path} & \textbf{Purpose}\\
\hline
\codepath{app/main.py} & Qt application entry point.\\
\codepath{app/dsp/} & Acquisition, tracking, bit sync, LNAV decoding, PVT, synthetic demos, I/O, compute backend selection, and display-oriented DSP helpers.\\
\codepath{app/gui/} & Main Qt window, worker plumbing, and tab widgets for each workflow stage.\\
\codepath{app/models/} & Shared dataclasses passed between GUI and DSP layers.\\
\codepath{app/tests/} & Unit, regression, equivalence, and smoke tests.\\
\codepath{tools/} & Local repository-management PowerShell helpers.\\
\end{longtable}

\section{Design Principles}

\begin{itemize}
  \item DSP logic lives in \codepath{app/dsp/}; GUI classes call it but do not hide signal-processing algorithms inside widgets.
  \item User-visible evidence stays PRN-specific: acquisition candidates, tracking histories, navigation words, and PVT inputs preserve the satellite identity that produced them.
  \item Large recordings are handled through both RAM-loaded arrays and bounded file windows. The reader in \codepath{app/dsp/io.py} is the central file-access layer.
  \item Optional acceleration is isolated. CPU execution remains the reference path; GPU-specific code is kept in focused backend helpers such as \codepath{app/dsp/tracking_gpu.py}.
  \item Real-capture diagnosis is first-class. Sample-rate surveys, search-center sweeps, repeated-segment acquisition evidence, and per-PRN interpretations are exposed rather than collapsed into one opaque score.
  \item PVT is offline and evidence-driven. It uses decoded LNAV ephemerides and testable pseudorange observations, without network assistance or map services.
\end{itemize}

\section{Runtime Dependencies}

The repository dependency file is intentionally small:

\begin{itemize}
  \item \texttt{numpy}: numerical arrays, FFTs, vectorized DSP, and tests.
  \item \texttt{pyqtgraph}: plotting widgets embedded in the Qt GUI.
  \item \texttt{PySide6}: Qt bindings for the desktop application.
\end{itemize}

CuPy/CUDA support is optional and is detected at runtime by \codepath{app/dsp/compute.py}. The application must remain runnable with:

\begin{lstlisting}[style=pythonapi]
python -m app.main
\end{lstlisting}

\section{DSP Package}

\subsection{\codeheading{app/dsp/acquisition.py}}

\textbf{Role.} Performs GPS L1 C/A acquisition over Doppler bins and code phase. It supports single-PRN acquisition, multi-PRN scans, search-center sweeps, and sample-rate surveys. The module emphasizes diagnosable repeated-segment evidence rather than only a single heatmap peak.

\textbf{Core data structures.}
\begin{itemize}
  \apiitem{AcquisitionConfig}{Compact settings for PRN, sample rate, Doppler search, code-period selection, workers, backend, and repeated-segment behavior.}
  \apiitem{\_AcquisitionPlan}{Precomputed vectors shared by segment searches.}
  \apiitem{\_SegmentSearchResult}{Best candidate plus heatmap for one searched recording segment.}
\end{itemize}

\textbf{Key entry points.}
\begin{itemize}
  \apiitem{acquire_signal(samples, config, progress_callback=None, log_callback=None)}{Runs one PRN acquisition and returns a full \texttt{AcquisitionResult}.}
  \apiitem{acquisition_from_session(samples, session, ...)}{Builds \texttt{AcquisitionConfig} from GUI/session settings.}
  \apiitem{scan_prns_from_session(samples, session, prns, ...)}{Acquires several PRNs and returns per-satellite results.}
  \apiitem{sweep_search_centers_from_session(samples, session, search_centers_hz, prns, ...)}{Ranks IF/search-center hypotheses.}
  \apiitem{survey_sample_rates(samples, session, sample_rates_hz, prns, ...)}{Ranks sample-rate hypotheses for ambiguous recordings.}
\end{itemize}

\textbf{Important helpers.}
\begin{itemize}
  \apisym{\_validate_acquisition_inputs}
  \apisym{\_prepare_acquisition}
  \apisym{\_select_ms_blocks}
  \apisym{\_select_segment_starts_ms}
  \apisym{\_build_heatmap_cpu}
  \apisym{\_build_heatmap_gpu}
  \apisym{\_top_heatmap_candidates}
  \apisym{\_cluster_segment_candidates}
  \apisym{acquisition_metric_is_strong}
  \apisym{acquisition_result_is_plausible}
  \apisym{acquisition_interpretation}
  \apisym{acquisition_rank_key}
\end{itemize}

\subsection{\codeheading{app/dsp/benchmark.py}}

\textbf{Role.} Benchmarks file I/O and DSP components so the GUI can report whether the current laptop is suitable for large offline recordings.

\textbf{Entry points.}
\begin{itemize}
  \apiitem{run_benchmark(session, progress_callback=None, log_callback=None)}{Measures representative I/O, FFT, acquisition-like, and tracking-like workload components.}
  \apiitem{\_benchmark_fft_core}{Runs the FFT core using the resolved compute backend.}
  \apiitem{\_component_result}{Normalizes each measured component into a display row.}
  \apiitem{\_system_info}{Collects lightweight system information.}
\end{itemize}

\subsection{\codeheading{app/dsp/bitsync.py}}

\textbf{Role.} Converts millisecond prompt integrations from tracking into GPS navigation-bit decisions.

\textbf{Entry points.}
\begin{itemize}
  \apiitem{form_navigation_bits(prompt_ms)}{Finds the best 20 ms boundary and forms hard bit decisions.}
  \apiitem{extract_navigation_bits(tracking)}{Applies carrier-aligned prompt extraction to a \texttt{TrackingState}.}
  \apiitem{\_carrier_aligned_prompt_ms}{Applies simple BPSK phase alignment before bit integration.}
\end{itemize}

\subsection{\codeheading{app/dsp/compute.py}}

\textbf{Role.} Centralizes CPU worker policy, optional CuPy/CUDA detection, and common parallel execution helpers.

\textbf{Data structures.}
\begin{itemize}
  \apiitem{GPUInfo}{Runtime status for optional GPU support.}
  \apiitem{ComputePlan}{Resolved backend and worker-count decision for an operation.}
  \apiitem{ProgressTracker}{Aggregates subtask progress into a monotonic overall percentage.}
\end{itemize}

\textbf{Entry points.}
\begin{itemize}
  \apiitem{detect_logical_cores}{Returns available logical CPU cores.}
  \apiitem{bootstrap_cuda_runtime_paths}{Adds pip-installed NVIDIA CUDA DLL paths to the current process environment when present.}
  \apiitem{get_cupy_module}{Returns CuPy when importable and usable.}
  \apiitem{detect_gpu_info}{Reports whether the optional GPU backend is available.}
  \apiitem{resolve_worker_count}{Chooses an executable CPU worker count from user preference and task count.}
  \apiitem{resolve_compute_plan}{Converts requested backend/workers into a concrete plan.}
  \apiitem{split_nested_worker_budget}{Divides workers across nested parallel loops.}
  \apiitem{parallel_ordered_map}{Runs independent tasks in parallel while preserving input order.}
\end{itemize}

\subsection{\codeheading{app/dsp/concept_lab.py}}

\textbf{Role.} Generates tiny deterministic GPS-like signals for the GUI concept lab so users can see carrier wipeoff, PRN separation, and correlation behavior without a real capture.

\textbf{Data structures.}
\begin{itemize}
  \apiitem{ConceptLabConfig}{Synthetic signal parameters.}
  \apiitem{ConceptLabResult}{Derived arrays displayed by the concept lab tab.}
\end{itemize}

\textbf{Entry points.}
\begin{itemize}
  \apiitem{generate_concept_lab_signal(config)}{Builds the deterministic teaching signal and correlation products.}
  \apisym{\_samples_per_ms}
  \apisym{\_phase_samples_to_chips}
  \apisym{\_code_for_phase}
  \apisym{\_correlation_profile}
  \apisym{\_component}
\end{itemize}

\subsection{\codeheading{app/dsp/demo.py}}

\textbf{Role.} Creates synthetic GPS-like baseband signals for demos and tests.

\textbf{Entry points.}
\begin{itemize}
  \apiitem{generate_demo_signal(sample_rate, duration_s, prn, doppler_hz, code_phase_samples, cn0_like_amplitude, noise_std)}{Generates a simplified PRN-coded BPSK signal with deterministic LNAV-like data.}
  \apiitem{\_build_demo_nav_bits}{Constructs deterministic navigation bits with an LNAV-style preamble and valid parity words.}
\end{itemize}

\subsection{\codeheading{app/dsp/ephemeris.py}}

\textbf{Role.} Decodes GPS LNAV ephemerides from valid subframes and computes satellite clock corrections and ECEF positions. The implementation stays close to the broadcast equations so the math remains teachable.

\textbf{Data structures.}
\begin{itemize}
  \apiitem{GpsEphemeris}{Broadcast LNAV ephemeris for one GPS PRN.}
  \apiitem{\_ClockFields}{Internal parsed subframe-1 clock fields.}
  \apiitem{\_OrbitPart2}{Internal parsed subframe-2 orbit fields.}
  \apiitem{\_OrbitPart3}{Internal parsed subframe-3 orbit fields.}
\end{itemize}

\textbf{Entry points.}
\begin{itemize}
  \apiitem{decode_ephemeris(prn, subframes)}{Builds one complete ephemeris from subframes 1, 2, and 3.}
  \apiitem{gps_time_delta_seconds(delta_s)}{Wraps GPS time differences to half-week range.}
  \apiitem{satellite_clock_correction_s(ephemeris, transmit_time_s)}{Computes satellite clock correction.}
  \apiitem{satellite_position_ecef_m(ephemeris, transmit_time_s)}{Computes satellite ECEF position at transmit time.}
  \apiitem{rotate_ecef_for_transit(ecef_m, travel_time_s)}{Rotates transmit-time ECEF coordinates into the receive-time frame.}
\end{itemize}

\subsection{\codeheading{app/dsp/gps_ca.py}}

\textbf{Role.} Generates and samples GPS L1 C/A PRN codes.

\textbf{Entry points.}
\begin{itemize}
  \apiitem{generate_ca_code(prn)}{Returns one 1023-chip PRN code as \texttt{+/-1} values.}
  \apiitem{sample_ca_code(prn, sample_rate, num_samples, code_phase_chips=0, code_rate_hz=1.023e6)}{Samples a local code replica at a requested sample rate.}
  \apiitem{code_phase_samples_to_chips(code_phase_samples, sample_rate)}{Converts a sample-index code phase into chips.}
\end{itemize}

\subsection{\codeheading{app/dsp/io.py}}

\textbf{Role.} Provides all raw IQ inspection and loading behavior for little-endian \texttt{complex64} recordings.

\textbf{Data structures.}
\begin{itemize}
  \apiitem{Complex64FileSource}{Memory-mapped, windowed reader for very large IQ files.}
\end{itemize}

\textbf{Entry points.}
\begin{itemize}
  \apiitem{common_sample_rate_hints(total_samples)}{Returns durations for common GNSS sample-rate hypotheses.}
  \apiitem{inspect_complex64_file(file_path, sample_rate, preview_samples)}{Returns metadata and preview statistics.}
  \apiitem{load_complex64_samples(file_path, start_sample, sample_count)}{Loads a bounded sample window.}
  \apiitem{load_complex64_samples_with_progress}{Loads a bounded window in chunks.}
  \apiitem{load_complex64_file}{Loads an entire recording into RAM.}
  \apiitem{load_complex64_file_with_progress}{Loads an entire recording into RAM with progress reporting.}
\end{itemize}

\subsection{\codeheading{app/dsp/navdecode.py}}

\textbf{Role.} Performs first-step GPS LNAV framing: preamble search, 30-bit word alignment, parity checks, conservative single-bit repair, TLM/HOW parsing, and GUI-ready subframe fields.

\textbf{Entry points.}
\begin{itemize}
  \apiitem{compute_lnav_parity(data_bits, d29_star, d30_star)}{Computes GPS LNAV parity bits \texttt{D25..D30}.}
  \apiitem{check_lnav_word(word_bits, previous_word=None)}{Checks parity for one 30-bit LNAV word.}
  \apiitem{extract_data_bits(word, previous_word=None)}{Recovers de-whitened 24-bit data payload.}
  \apiitem{parse_tlm(word, previous_word=None)}{Decodes teaching/display TLM fields.}
  \apiitem{parse_how(word, previous_word=None)}{Decodes TOW count and subframe ID.}
  \apiitem{classify_subframe(subframe_id)}{Returns a didactic label for subframe type.}
  \apiitem{maybe_correct_word(word_bits, previous_word, confidences)}{Attempts one unambiguous low-confidence bit repair.}
  \apiitem{build_subframes(values, preamble_indices, inverted_flag, confidences)}{Groups bit streams into candidate LNAV subframes.}
  \apiitem{decode_navigation_bits(bits)}{Searches for preambles and decodes LNAV framing.}
  \apiitem{decode_navigation_from_tracking(tracking, bit_source, progress_callback=None, log_callback=None)}{Runs bit extraction and LNAV decoding from tracking output.}
\end{itemize}

\subsection{\codeheading{app/dsp/pvt.py}}

\textbf{Role.} Implements receiver position, velocity/time helper routines and the offline PVT builder that consumes decoded navigation evidence.

\textbf{Data structures.}
\begin{itemize}
  \apiitem{PositionSolution}{Least-squares receiver ECEF/LLA position and clock bias.}
  \apiitem{PseudorangeObservation}{One observation formed from valid LNAV and tracking timing evidence.}
  \apiitem{PVTComputationResult}{End-to-end PVT result and supporting evidence.}
  \apiitem{\_SubframeTiming}{Internal receive-time evidence for a decoded subframe.}
  \apiitem{\_SubframeArrival}{Internal decoded subframe arrival tied to file time.}
  \apiitem{\_ClockCorrectedArrival}{Internal arrival after satellite clock correction.}
\end{itemize}

\textbf{Entry points.}
\begin{itemize}
  \apiitem{lla_to_ecef(latitude_deg, longitude_deg, altitude_m)}{Converts WGS-84 geodetic coordinates to ECEF metres.}
  \apiitem{ecef_to_lla(ecef_m)}{Converts ECEF metres to WGS-84 latitude, longitude, and altitude.}
  \apiitem{solve_position_from_pseudoranges(satellite_positions_m, pseudoranges_m)}{Solves receiver ECEF position and clock bias from at least four satellites.}
  \apiitem{expand_gps_week(week_number_mod1024, reference)}{Expands a 10-bit LNAV week number near a reference date.}
  \apiitem{gps_utc_datetime(gps_week, time_of_week_s, leap_seconds=18)}{Converts GPS week/TOW to UTC.}
  \apiitem{compute_pvt_from_navigation(tracking_by_prn, bit_results_by_prn, nav_results_by_prn, ...)}{Builds ephemerides, pseudoranges, receiver position, and GPS/UTC time.}
\end{itemize}

\subsection{\codeheading{app/dsp/pvt_pipeline.py}}

\textbf{Role.} Wires existing DSP stages for the PVT tab: acquire several PRNs, track them, decode LNAV, then call the PVT computation layer.

\textbf{Data structures and entry points.}
\begin{itemize}
  \apiitem{PVTPipelineResult}{Contains outputs from the automated PVT teaching pipeline.}
  \apiitem{run_pvt_pipeline(file_path, session)}{Runs acquisition, tracking, LNAV decode, and PVT solve on one file window.}
  \apiitem{\_pvt_session}{Builds a session variant for pipeline defaults.}
\end{itemize}

\subsection{\codeheading{app/dsp/tracking.py}}

\textbf{Role.} Implements the educational carrier/code tracking loop after acquisition. It supports RAM arrays and streamed file windows, CPU and optional GPU vector math, and PRN-specific diagnostics.

\textbf{Internal data structures.}
\begin{itemize}
  \apiitem{\_TrackingLoop}{Mutable carrier/code state.}
  \apiitem{\_CorrelatorOutput}{Prompt, early, and late outputs for one millisecond.}
  \apiitem{\_TrackingMath}{Small NumPy/CuPy adapter.}
  \apiitem{\_TrackingHistory}{Preallocated arrays filled by the loop.}
\end{itemize}

\textbf{Public entry points.}
\begin{itemize}
  \apiitem{track_signal(samples, session, acquisition, source_start_sample=0, ...)}{Tracks one PRN from an in-memory sample array.}
  \apiitem{track_file(file_path, start_sample, session, acquisition, ...)}{Tracks one PRN directly from a large file without loading the full recording.}
\end{itemize}

\textbf{Important helpers.}
\begin{itemize}
  \apisym{\_validate_tracking_inputs}
  \apisym{\_initial_tracking_loop}
  \apisym{\_tracking_math_for_backend}
  \apisym{\_sample_ca_code_for_loop}
  \apisym{\_wipe_carrier}
  \apisym{\_correlate_prompt_early_late}
  \apisym{\_correlate_block_gpu}
  \apisym{\_loop_discriminators}
  \apisym{\_advance_tracking_loop}
  \apisym{\_detect_lock}
  \apisym{\_build_tracking_state}
  \apisym{\_run_tracking_loop}
  \apisym{\_iter_sample_blocks}
\end{itemize}

\subsection{\codeheading{app/dsp/tracking_gpu.py}}

\textbf{Role.} Contains the optional CuPy raw kernel for tracking correlators. Keeping this file separate preserves a readable CPU-oriented teaching path in \codepath{tracking.py}.

\textbf{Entry point.}
\begin{itemize}
  \apiitem{get_tracking_correlator_kernel()}{Compiles and returns the optional CuPy E/P/L correlator kernel.}
\end{itemize}

\subsection{\codeheading{app/dsp/utils.py}}

\textbf{Role.} Provides display-oriented DSP helpers for decimation, smoothing, windows, spectrum, waterfall, and bit formatting.

\textbf{Entry points.}
\begin{itemize}
  \apiitem{decimate_for_display(values, max_points)}{Returns a compact copy suitable for plotting.}
  \apiitem{moving_average(values, window)}{Smooths display traces.}
  \apiitem{get_window(window_name, length)}{Returns a standard FFT window.}
  \apiitem{compute_spectrum(samples, sample_rate, fft_size, window_name, average_count, session=None)}{Computes averaged spectrum data.}
  \apiitem{compute_waterfall(samples, sample_rate, fft_size, step, window_name, max_rows, session=None)}{Builds a compact waterfall matrix.}
  \apiitem{bits_to_str(bits)}{Formats bit arrays as a compact string.}
\end{itemize}

\section{Models}

\subsection{\codeheading{app/models/session.py}}

\textbf{Role.} Defines the shared dataclasses that form the contract between GUI tabs, worker threads, DSP modules, and tests.

\textbf{Session and file state.}
\begin{itemize}
  \apiitem{SessionConfig}{User-selected analysis settings such as file path, sample rate, PRN, acquisition/tracking limits, backend, preload/windowing, and PVT settings.}
  \apiitem{FileMetadata}{Raw IQ file description, sample count, duration, preview statistics, and rate-hypothesis hints.}
\end{itemize}

\textbf{Acquisition state.}
\begin{itemize}
  \apiitem{AcquisitionCandidate}{One heatmap peak candidate with PRN, Doppler, code phase, metric, and segment evidence.}
  \apiitem{AcquisitionResult}{Full acquisition output for one PRN and search-center hypothesis.}
  \apiitem{SearchCenterSweepEntry}{One ranked search-center result.}
  \apiitem{SearchCenterSweepResult}{A collection of search-center sweep entries.}
  \apiitem{SampleRateSurveyEntry}{One ranked sample-rate hypothesis.}
  \apiitem{SampleRateSurveyResult}{A collection of sample-rate survey entries.}
\end{itemize}

\textbf{Tracking and navigation state.}
\begin{itemize}
  \apiitem{TrackingState}{Per-millisecond prompt/early/late, carrier/code estimates, lock evidence, and IQ previews.}
  \apiitem{BitDecisionResult}{Navigation-bit decisions, confidences, polarity, and bit-boundary evidence.}
  \apiitem{NavigationField}{One decoded or coarse LNAV display field.}
  \apiitem{NavigationWord}{One 30-bit LNAV word with parity and repair metadata.}
  \apiitem{NavigationSubframe}{One 10-word LNAV subframe view.}
  \apiitem{NavigationDecodeResult}{Preamble, inversion, subframe, and field results from nav decoding.}
\end{itemize}

\textbf{Benchmark and demo state.}
\begin{itemize}
  \apiitem{BenchmarkComponentResult}{One measured subsystem row.}
  \apiitem{BenchmarkResult}{Aggregate benchmark status and component rows.}
  \apiitem{DemoSignalResult}{Synthetic signal plus reference truth values.}
\end{itemize}

\section{GUI Package}

\subsection{\codeheading{app/main.py}}

\textbf{Role.} Creates the Qt application, instantiates \texttt{MainWindow}, and starts the event loop through \texttt{main()}.

\subsection{\codeheading{app/gui/workers.py}}

\textbf{Role.} Keeps long-running DSP actions off the UI thread.

\begin{itemize}
  \apiitem{WorkerSignals}{Qt signals for result, error, progress, log, and finished notifications.}
  \apiitem{Worker}{Generic \texttt{QRunnable} wrapper around a callable.}
\end{itemize}

\subsection{\codeheading{app/gui/main_window.py}}

\textbf{Role.} Top-level Qt controller. It owns the session state, tab instances, loaded sample windows, background workers, and transitions between acquisition, tracking, navigation decode, and PVT.

\textbf{Major responsibilities.}
\begin{itemize}
  \item File selection, metadata inspection, RAM/windowed loading decisions, and display-sample previews.
  \item Session synchronization from GUI controls.
  \item Stale-result protection through processing-context signatures.
  \item Acquisition, all-PRN scan, search-center sweep, and sample-rate survey orchestration.
  \item Tracking orchestration for selected PRNs.
  \item Navigation decode orchestration from tracking outputs.
  \item PVT solve and automated PVT pipeline orchestration.
  \item Demo signal generation and benchmark dispatch.
\end{itemize}

\textbf{Important methods.}
\begin{itemize}
  \apisym{inspect_file}
  \apisym{ensure_samples}
  \apisym{load_display_samples}
  \apisym{load_acquisition_samples}
  \apisym{current_window_samples}
  \apisym{start_acquisition}
  \apisym{scan_all_prns}
  \apisym{start_search_center_sweep}
  \apisym{start_sample_rate_survey}
  \apisym{start_tracking}
  \apisym{decode_navigation}
  \apisym{solve_pvt_from_decoded}
  \apisym{start_pvt_pipeline}
  \apisym{generate_demo}
  \apisym{start_benchmark}
\end{itemize}

\subsection{GUI Tabs}

\begin{longtable}{P{0.34\linewidth}L{0.58\linewidth}}
\textbf{Path} & \textbf{Class and responsibility}\\
\hline
\codepath{app/gui/tabs/session_tab.py} & \texttt{SessionTab}: file loading, session parameters, RAM/backend status, progress, and log output.\\
\codepath{app/gui/tabs/raw_signal_tab.py} & \texttt{RawSignalTab}: time-domain views of raw IQ samples.\\
\codepath{app/gui/tabs/iq_tab.py} & \texttt{IQPlaneTab}: IQ scatter and phasor diagnostics for selected sources.\\
\codepath{app/gui/tabs/spectrum_tab.py} & \texttt{SpectrumTab}: averaged FFT spectrum and waterfall display.\\
\codepath{app/gui/tabs/acquisition_tab.py} & \texttt{AcquisitionTab}: PRN list parsing, acquisition heatmap, Doppler/code-phase slices, PRN/Doppler overview, search-center sweep, and sample-rate survey displays.\\
\codepath{app/gui/tabs/tracking_tab.py} & \texttt{TrackingTab}: lock, prompt/early/late, DLL/PLL, code/carrier estimates, and IQ-stage diagnostics.\\
\codepath{app/gui/tabs/navigation_tab.py} & \texttt{NavigationTab}: bit extraction, LNAV words, subframes, fields, and PRN-specific navigation results.\\
\codepath{app/gui/tabs/pvt_tab.py} & \texttt{PVTTab}: PVT pipeline settings, receiver solution, ephemeris table, and observation table.\\
\codepath{app/gui/tabs/benchmark_tab.py} & \texttt{BenchmarkTab}: laptop suitability and component throughput results.\\
\codepath{app/gui/tabs/concept_lab_tab.py} & \texttt{ConceptLabTab}: synthetic signal concept demonstrations and plots.\\
\codepath{app/gui/tabs/learning_tab.py} & \texttt{LearningTab}: compact explanation of the current receiver pipeline state.\\
\end{longtable}

\subsection{\codeheading{app/gui/tabs/acquisition_tab.py}}

\textbf{Notable methods.}
\begin{itemize}
  \apisym{parse_prn_list}
  \apisym{selected_scan_prns}
  \apisym{build_prn_doppler_overview}
  \apisym{threshold_prn_doppler_overview}
  \apisym{sparse_prn_axis_ticks}
  \apisym{overview_marker_rows}
  \apisym{best_heatmap_indices}
  \apisym{codephase_slice}
  \apisym{doppler_slice}
  \apisym{update_result}
  \apisym{update_sweep_results}
  \apisym{update_sample_rate_survey}
  \apisym{build_search_centers}
\end{itemize}

\subsection{\codeheading{app/gui/tabs/navigation_tab.py}}

\textbf{Notable methods.}
\begin{itemize}
  \apisym{set_available_prns}
  \apisym{bit_source_mode}
  \apisym{update_results}
  \apisym{\_refresh_decoded_fields_table}
  \apisym{\_format_almanac_ephemeris}
\end{itemize}

\subsection{\codeheading{app/gui/tabs/pvt_tab.py}}

\textbf{Notable methods.}
\begin{itemize}
  \apisym{pipeline_settings}
  \apisym{clear_result}
  \apisym{update_result}
  \apisym{\_update_ephemeris_table}
  \apisym{\_update_observation_table}
\end{itemize}

\section{Tests}

The tests document expected behavior and guard against regressions in DSP math, large-file behavior, GUI source switching, optional acceleration, and repository policy.

\begin{longtable}{P{0.36\linewidth}L{0.56\linewidth}}
\textbf{Test file} & \textbf{Coverage}\\
\hline
\codepath{app/tests/conftest.py} & Pytest import-path and Qt defaults.\\
\codepath{app/tests/test_acceleration_paths.py} & CPU multi-worker equivalence, optional GPU equivalence, and GPU-to-CPU fallback behavior.\\
\codepath{app/tests/test_acquisition_confidence.py} & Acquisition plausibility, ranking, and segment-consistency scoring.\\
\codepath{app/tests/test_acquisition_demo.py} & Synthetic acquisition truth, PRN identity, and invalid-window validation.\\
\codepath{app/tests/test_benchmark.py} & Benchmark result component sanity.\\
\codepath{app/tests/test_compute.py} & Worker-count policy and GPU detection edge cases.\\
\codepath{app/tests/test_concept_lab.py} & Concept-lab correlation peaks and PRN separability.\\
\codepath{app/tests/test_gps_ca.py} & C/A code length, values, reference starts, PRN taps, caching, and sampling validation.\\
\codepath{app/tests/test_gui_smoke.py} & Main window smoke tests plus acquisition, stale-context, source, and control-regression checks.\\
\codepath{app/tests/test_io.py} & Complex64 inspection, loading, sample-rate hints, and window capping.\\
\codepath{app/tests/test_large_file_io.py} & Windowed block reading from large-file source abstraction.\\
\codepath{app/tests/test_pvt.py} & PVT solver, GPS week expansion, UTC conversion, and tracked code-phase timing.\\
\codepath{app/tests/test_real_capture_helpers.py} & Real-capture diagnosis helpers: segment consistency and sample-rate survey.\\
\codepath{app/tests/test_repo_policy.py} & Repository delivery and sync policy visibility.\\
\codepath{app/tests/test_search_center_sweep.py} & Search-center sweep ranking and worker-budget splitting.\\
\codepath{app/tests/test_source_switching.py} & Demo-to-real-file source switching.\\
\codepath{app/tests/test_tracking_bits_nav.py} & LNAV parity, repair, subframe grouping, tracking validation, bit extraction, and nav pipeline progress.\\
\end{longtable}

\section{Repository Tools}

\subsection{\codeheading{tools/check_git_sync.ps1}}

\textbf{Role.} Quick pre-push verification helper for repository sync policy. It is intended to be run before sharing work when checking branch/remote state is useful.

\subsection{\codeheading{tools/install_git_hooks.ps1}}

\textbf{Role.} Installs local Git hooks from \codepath{.githooks/}. The repository currently includes a pre-push hook location.

\subsection{\codeheading{.githooks/pre-push}}

\textbf{Role.} Enforces the repository sync convention at push time. When pushing to \texttt{origin}, it verifies that \texttt{origin} has at least two push URLs configured so GitHub and GitLab can receive the same commit.

\section{Package Initializers}

\begin{longtable}{P{0.34\linewidth}L{0.58\linewidth}}
\textbf{File} & \textbf{Description}\\
\hline
\codepath{app/__init__.py} & Declares the GPS L1 C/A offline decoder package.\\
\codepath{app/dsp/__init__.py} & Declares the DSP package.\\
\codepath{app/gui/__init__.py} & Declares the Qt GUI package.\\
\codepath{app/gui/tabs/__init__.py} & Declares GUI tab widgets package.\\
\codepath{app/models/__init__.py} & Declares shared dataclasses package.\\
\end{longtable}

\section{End-to-End Data Flow}

\begin{enumerate}
  \item \texttt{SessionTab} and \texttt{MainWindow} collect a \texttt{SessionConfig}.
  \item \codepath{app/dsp/io.py} inspects or loads a recording window.
  \item \codepath{app/dsp/acquisition.py} searches PRNs and produces \texttt{AcquisitionResult} objects.
  \item \codepath{app/dsp/tracking.py} uses acquisition peaks to produce \texttt{TrackingState} histories.
  \item \codepath{app/dsp/bitsync.py} forms \texttt{BitDecisionResult} data from prompt integrations.
  \item \codepath{app/dsp/navdecode.py} builds \texttt{NavigationDecodeResult} subframes and fields.
  \item \codepath{app/dsp/ephemeris.py} decodes broadcast ephemerides from valid subframes.
  \item \codepath{app/dsp/pvt.py} forms observations and solves receiver position/time.
  \item GUI tabs render the evidence at each stage while keeping PRNs visually separated.
\end{enumerate}

\section{Build Notes for This Document}

This file is ordinary LaTeX and can be compiled with a standard distribution:

\begin{lstlisting}[style=pythonapi]
pdflatex docs/code_documentation.tex
\end{lstlisting}

No generated PDF is committed by default. The source file is the repository artifact; generated binaries should remain local unless explicitly needed for a release bundle.

\end{document}
